\documentclass[12pt]{article}
\usepackage{text_style}
\title{Forward and inverse problems in model generation}
\author{IA pour la Science}
\date{7 May 2026}
\begin{document}
\maketitle

Two aspects of the SCM inference were discussed the last week:
\begin{enumerate}
\item The inference avoids the  Integer Liner Programming problem by replacing the probabilities of edges
\[\bGam \in [0,1]^{n\times n} \mapsto \bZ \in \{0,1\}^{n\times n}.\]
for the adjacency of edges. \emph{NB: it is solved below.}
\item  The list of variables is not fixed.
For the set of observable variables $X_1,\dots,X_n$
the model has to forecast the list of variables $X_1,\dots,X_n, \dots, X_{n^\prime}$.
By default, the value $n^\prime=0$. \emph{NB: it is under consideration.}
\end{enumerate}

The foundation model includes the elements of SCM selection:
\begin{enumerate}[1)]
\item maximizing likelihood of the DAG at each step,
\item do-intervention actions each step on the sampled node.
\end{enumerate}

\subsection*{The foundation model for the SCM inference}
The main type of the foundation model carries the Riemannian diffusion process.

The diffusion process involves the step~$t$ and two distributions to approximate the distribution of 
the adjacency matrix~$\bZ$.
Use the vector form~$\bz = \text{vec}(\bZ)$.
Denote the training distribution~$q(\bz)$ and the sampling distribution~$p(\bz)$.
During the model training one has to minimize the distance between two distributions at each step. 

The diffusion process invokes the local structural modeling:
\[
\bZ \mapsto f_g \mapsto \eta \mapsto \bw,
\]
with $\eta$ is the optimization algorithm that delivers the optimal parameters to the SCM at each step of the diffusion process.

The \emph{fine-tuning} algorithm~$\eta$ optimizes the parameters according to the diagram

% Optimization and inference by an SCM.
%{\color{red}\footnote*{\color{gray}These notations are equivalent: $x_i, x(i), i\to x$.}}
\begin{equation*}
\xymatrix{
\eta(-\lambda^2\nabla_{\textbf{w}}\loss+\textbf{w}^\prime)\ar[r]^{\qquad \eta} & \hat{\bw}, {\bx} \ar[r]^{f_G} &  \hat{\bx}\ar[r] & \loss(\mathbf{\textbf{w}\mid\bx,\hat{\bx}}) \ar@/_3pc/[lll]\\
%1 & 2 &  & y\ar[ur] & 5 \\
(t,s)\ar[r] & i\ar[r] & \bx\ar[ul]\ar[ur] & }
\end{equation*}
This diagram reads: the optimization algorithm $\eta$ estimates the SCM parameters~$\bw$ by the gradient of its loss and the previous parameters to solve the problem 
\begin{equation*}
\hat{\textbf{w}}= \arg\min \loss(\bw) \quad\text{given}\quad f_G(\hat{\bw},\bx)\mapsto \hat{\bx}.
\end{equation*}
The algorithm $\eta$ uses some initial value~$\bw_0$ and a stop-criterion for its iterations.

The 

-
\end{document}








\subsection*{The empirical distribution of the observations}
$P(X_1,\dots,X_n)$ is the distribution of the observable data.
Its approximation is given by the Gaussian Process Regression model
\[
p(\bx \mid f_\text{gp})
\]
on the dataset~
\[\{\bx_i\in\mathbb{R}^n\}.\]
It is indexed by some function of the time and space $(t,s)\mapsto i$.

\subsection*{A model from the family of universal approximators}
For a sample $\bx$ and its reconstruction $\hat\bx$ the likelihood function provided a universal model~$f_\text{nn}(\bth,\bx)$ is
\[p(\bx\mid f_\text{nn}, \bth).\]
Without the a priori information about the model parameters~$\bth$ the corresponding loss function is
\[
\loss(\bth) = -\log p(\bx\mid f_\text{nn}, \bth) = (\hat{\bx} - \bx)\T \mathbf{B}^{-1}(\hat{\bx} - \bx).% + (\bth - \boldsymbol{\mu})\T \mathbf{A}^{-1}(\bth - \boldsymbol{\mu}),
    \]
The matrix $\mathbf{B}$ is the covariance of the reconstruction error.
This matrix in some statements depends on the index~$i$.
This model as a parametric one has a structure to make a distillation to the causal graph
\[
f_\text{nn}(\bth, \bx)  \overset{\text{extract}}{\longrightarrow} \bZ_\text{nn} \overset{\text{dist}}{\longrightarrow}  \bZ.
\]

\subsection*{An SCM we have to forecast}
A structural causal model
\[
    \hat{\bx} = f_G(\bw,\bx) \in\mathcal{F}_\text{scn}
\]
has a structure~$\bZ$ and its probabilistic representation $\bGam$.
The elements of the matrix~$\bGam$ are the probabilities of the causal relationships between the elements of the model structure~$f_G$.
The matrix~$\bZ$ is the adjacency matrix of the causal graph, the binary representation of~$\bGam$.

\medskip
Note: \emph{we suppose that SCM does not fit the whole $(s,t)$-data but a local part of it.
So we also call SCM a local model~$f_G$, in contrast to the universal model~$f_\text{nn}$.}

\subsection*{A foundational model to forecast the SCM}
There given pairs a sample $\bx_i$ and the corresponding SCM~$f_G$ with its structure $\{\bZ_i\}$.
Call the set of pairs
\[(\bx_i, \bZ_i), \quad i\in\cB
\]
a collection of the SCM forecasting.
This collection $\cB$ shows one case of mathematical modeling that results in an interpretable model~$f_G$.
This model fits samples $\bx_i$ from some local data~$i\in\cB$.

Call a model that forecasts an SCM $\mathcal{F}$ a \emph{foundational model}
\[
f_\text{ai}: \bx_\cB \mapsto f_G \in \mathcal{F}.
\]
The SCM is a parametric model
\[
    f_G = f_G (\bw,\bx) \quad\text{with the parameters}\quad \hat{\bw} = \arg\min_{\bw} \loss(\bw).
\]
The foundational model returns not the SCM~$f_G(\hat{\bw},\bx)$ itself,
but only a probabilities~$\bGam$ on the structure~$\bZ$.
So we rewrite the previous definition as
\[
f_\text{ai}: \bx_\cB \mapsto \bGam.
\]

\medskip
Note: \emph{Since a case of mathematical modeling, call it $\cB$ has a history of requests to the foundational model,
    the next information is available for it:
\begin{enumerate}[1)]
    \item the optimal parameters $\bw$ along with $\bZ, \bGam$ for the SCM $f_G$,
    \item the prior probabilities of the model parameters $p(f_G)$ or all models in the family $\mathcal{F}$,
    \item the other information about the case of mathematical modeling $\cB$ that is relevant for the forecasting of the SCM.
\end{enumerate}}

\subsection*{The Integer Linear Programming for the causal graph selection}
The matrix $\bGam$ comes along with the descriptions of its nodes $X_1,\dots,X_n$, and the probabilities.
The SCM $f_G$ is defined by the graph $G=(\bZ,X)$ by solving
\[
  \text{ILP}(\bGam,P_X) \to \bZ \to f_G.
\]

\subsection*{The evidence of the SCM}
A probabilistic description~$\bGam$ of a model $f_G$ from the class of \emph{admissible models} $\mathcal{F}$ is generated from a parametric distribution
\[
    \bGam \sim \text{Dir}(\bpi) \quad\text{generated by MCMC}\quad f_\text{mh}(\bpi),
\]
here the lowercase letters mean Metropolis-Hastings Algorithm for MCMC.
(A simple genetic algorithm may bring some decent results, too).
This distribution provides a probabilistic space for the similar models.
See the document on \emph{Multimodeling}.

The model $f_G$ and the corresponding graph has its likelihood
\[
    p(\bx|f_G).
\]
See the document on \emph{Causal graph likelihood}.

\medskip
Note: \emph{Two-level bayesian inference introduces the probability of the SCM parameters on the first level.
And the probability of the model on the second level is:
    \begin{enumerate}[1)]
    \item the prior of the SCM is delivered by the generative model $f_\text{mh}$
    \item the likelihood of the SCM is delivered by the forecasting model $f_\text{ai}$.
\end{enumerate}}

\subsection*{The criterion of causality}
The parameters of the model $f_G$ has their covariance matrix.
The prior of the model parameters is
\[
    - \ln p(\bw|f_g) = (\bw - \bmu)^{\mathsf{T}} \bA^{-1} (\bw - \bmu).
\]
The algorithm of the Bayesian hyperparameters $(\bmu, \bA)$ estimation
(and their connection between the causality)
is described in the document on \emph{parameter analysis}.

\subsection*{A model for the intervention}
See the document on the \emph{intervention strategy}.

\end{document}
